% ch3_f 第3章 §3.6尾–§3.7 (书页 123–133 / p138–p148)
% 边界备注：
%   起点JOIN——书页 123 顶部 "of the density correlation function:" 是书页 122 末句
%   "Let us consider the spectral expansion of the density correlation function:" 的后半句；
%   完整引导句（"让我们考虑密度关联函数的谱展开："）应由上一块的 E-TAIL 收入，
%   本块正文以 %JOIN% 起头，直接从其后的谱展开公式开始。
%   终点——书页 133 止于式 v_c = Min(ε(k)/|k|)，句意完结；书页 134 顶部为图 3.15 与
%   新段落 "It is interesting to see that ..."，属下一块，本块自然结束，无 TAIL。
%   本块范围内无插图（图 3.12–3.14 在书页 121–122，图 3.15 在书页 134），
%   正文引用图 3.1、图 3.15(a)(b)；脚注共 1 处（原书脚注 24）。

%JOIN%
\begin{align*}
\langle\rho(t,x)\rho(0,0)\rangle&=\frac{\langle0|\rho(x)U(t,0)\rho(0)|0\rangle}{\langle0|U(t,0)|0\rangle}\\
&=\sum_{n,k}\langle0|\rho(x)|n,k\rangle\langle n,k|\rho(0)|0\rangle\,\mathrm{e}^{\mathrm{i}kx-\mathrm{i}\epsilon_{n,k}t}
\end{align*}
其中态 $|n,k\rangle$ 具有能量 $\epsilon_{n,k}$ 和动量 $k$。我们还假定了基态 $|0\rangle$ 的能量为零。由于低能态只出现在 $K_i$ 附近，在低能下我们可以把上述求和重新分组如下：
\begin{equation*}
\langle\rho(t,x)\rho(0,0)\rangle=\sum_{n,|\delta k|\ll K_1,i}\langle0|\rho(x)|n,\delta k,i\rangle\langle n,\delta k,i|\rho(0)|0\rangle\,\mathrm{e}^{\mathrm{i}(\delta k+K_i)x-\mathrm{i}\epsilon_{n,\delta k,i}t}
\end{equation*}
其中态 $|n,\delta k,i\rangle$ 具有能量 $\epsilon_{n,\delta k,i}$ 和动量 $\delta k+K_i$。我们看到，大动量处的密度关联，比如 $k\sim K_i$ 处的密度关联，是由矩阵元 $\langle n,\delta k,i|\rho(0)|0\rangle$ 产生的。我们知道，涡旋算符 $O_v^i$（当 $i<0$ 时 $O_v^i\equiv(O_v^\dagger)^{-i}$）把 $k=0$ 附近的低能态映射到 $k=K_i$ 附近的低能态。于是这里我们作如下大胆的假设：
\begin{equation*}
\langle n,\delta k,i|\rho(t,x)|0\rangle=C_i\langle n,\delta k,i|O_v^i(t,x)|0\rangle
\end{equation*}
低能密度算符现在可以推广为
\begin{equation*}
\rho(t,x)=\rho_0-\chi\dot{\theta}(t,x)+\sum_n C_nO_v^n
\end{equation*}
旧结果 $\rho(t,x)=\rho_0-\chi\dot{\theta}(t,x)$ 只描写低能长波长密度涨落，新结果则描写低能全波长密度涨落。我们发现关联函数具有如下形式：
\begin{align*}
\langle\rho(t,x)\rho(0,0)\rangle
&=\rho_0^2-\frac{\chi v}{4\pi}\left(\frac{1}{(x-vt)^2}+\frac{1}{(x+vt)^2}\right)+\sum_n|C_n|^2\mathrm{e}^{\mathrm{i}K_nx}\left(\frac{l^2}{x^2-v^2t^2}\right)^{n^2\pi\chi v}\\
&=\rho_0^2-\left(\frac{\chi v/4\pi}{(x-vt)^2}+\frac{\chi v/4\pi}{(x+vt)^2}\right)+\sum_n|C_n|^2\mathrm{e}^{\mathrm{i}K_nx}\left(\frac{l^2\mathrm{e}^{-\mathrm{i}\pi\Theta(v^2t^2-x^2)}}{|x^2-v^2t^2|}\right)^{n^2\pi\chi v}\tag{3.6.2}
\end{align*}
其中 $l$ 是短距离截断。注意，虚时中的关联 $\langle O_v^\dagger(x,t)O_v(0,0)\rangle$ 由 $\mathrm{e}^{-V}$ 给出，其中 $V$ 是涡旋/反涡旋对之间的势。实时关联则由解析延拓得到。

\subsection*{问题 3.6.1}

证明式 (3.6.2)。

\subsection*{问题 3.6.2}

有限动量下的响应率。
\begin{enumerate}
\item 设玻色子感受到一个弱势 $V(x)$，该势引起密度改变 $\delta\rho$。试把 $\delta\rho_k=\chi(k)V_k$ 中的有限动量响应率 $\chi(k)$ 用频率--动量空间中的密度关联函数表示出来。
\item 求出使响应率 $\chi(K_n)$ 发散的 $\chi$ 与 $v$ 的取值。（提示：先在虚时中做计算可能会容易一些。）
\item 如果我们开启一个周期 $a=\rho^{-1}/n$（$n$ 为整数）的弱周期势，那么对哪些 $\chi$ 与 $v$ 的取值，玻色子形成莫特绝缘体？
\end{enumerate}

\section{超流与超导}

\subsection{与规范场的耦合和守恒流}

\begin{keypoints}
\item 具有整体 $U(1)$ 对称性的理论包含一个守恒荷。
\item 具有整体 $U(1)$ 对称性的理论可以与 $U(1)$ 规范场耦合。电磁矢量势就是一个 $U(1)$ 规范场。
\end{keypoints}

带电玻色子系统会与电磁规范场耦合。当电磁场不为零时，带电玻色子系统的拉格朗日量需要修改。这里的问题是：我们如何找到修改后的拉格朗日量？我们知道，玻色子拉格朗日量
\begin{equation*}
L(\varphi)=\mathrm{i}\frac{1}{2}(\varphi^*\partial_t\varphi-\varphi\partial_t\varphi^*)-\frac{1}{2m}\partial_{\pmb{x}}\varphi^*\partial_{\pmb{x}}\varphi+\mu|\varphi|^2-\frac{V_0}{2}|\varphi|^4 \tag{3.7.1}
\end{equation*}
在整体 $U(1)$ 变换
\begin{equation*}
\varphi\rightarrow\mathrm{e}^{\mathrm{i}f}\varphi
\end{equation*}
下不变，即 $L(\mathrm{e}^{\mathrm{i}f}\varphi)=L(\varphi)$。但它在一个局域 $U(1)$ 变换
\begin{equation*}
\varphi(\pmb{x},t)\rightarrow\mathrm{e}^{\mathrm{i}f(\pmb{x},t)}\varphi(\pmb{x},t)
\end{equation*}
下并非不变。我们有
\begin{align*}
L(\mathrm{e}^{\mathrm{i}f(\pmb{x},t)}\varphi)&=\mathrm{i}\frac{1}{2}\big(\varphi^*(\partial_0+\mathrm{i}\partial_0f)\varphi-\varphi(\partial_0-\mathrm{i}\partial_0f)\varphi^*\big)\\
&\quad-\frac{1}{2m}\big|(\partial_i+\mathrm{i}\partial_if)\varphi\big|^2+\mu|\varphi|^2-\frac{V_0}{2}|\varphi|^4 \tag{3.7.2}
\end{align*}
其中下标 0 标记时间方向，下标 $i=1,\dots,d$ 标记空间方向。我们将用希腊字母 $\mu,\nu$ 等标记时空方向。例如，$x^\mu$ 代表时空坐标。玻色子与电磁规范场的耦合现在只需把 $\partial_\mu f$ 换成 $A_\mu$ 便可得到：
\begin{align*}
L(\varphi,A_\mu)&=\mathrm{i}\frac{1}{2}\big(\varphi^*(\partial_0+\mathrm{i}A_0)\varphi-\varphi(\partial_0-\mathrm{i}A_0)\varphi^*\big)\\
&\quad-\frac{1}{2m}\big|(\partial_i+\mathrm{i}A_i)\varphi\big|^2+\mu|\varphi|^2-\frac{V_0}{2}|\varphi|^4 \tag{3.7.3}
\end{align*}
所得的拉格朗日量有一个很好的性质：它是规范不变的：
\begin{gather*}
\varphi\rightarrow\tilde{\varphi}=\mathrm{e}^{\mathrm{i}f(\pmb{x},t)}\varphi\\
A_\mu\rightarrow\tilde{A}_\mu=A_\mu-\partial_\mu f\\
L(\varphi,A_\mu)\rightarrow L(\tilde{\varphi},\tilde{A}_\mu)=L(\varphi,A_\mu) \tag{3.7.4}
\end{gather*}
上述构造并不包含多少物理，它只是得到规范不变拉格朗日量的一个技巧。这一技巧给出最简单的拉格朗日量，称为最小耦合（minimal coupling）。我们也可以写出别的规范不变拉格朗日量，例如把 $\partial_\mu f$ 换成 $A_\mu+g\partial_\nu F_{\nu\mu}$，其中 $F_{\mu\nu}=\partial_\mu A_\nu-\partial_\nu A_\mu$ 是 $A_\mu$ 的场强。注意 $F_{\mu\nu}$ 在规范变换下不变。

式 (3.7.3) 只描写荷场 $\varphi$ 与规范场 $A_\mu$ 之间的耦合。同时描写 $\varphi$ 与 $A_\mu$ 二者动力学的完整规范不变拉格朗日量为
\begin{align*}
L(\varphi,A_\mu)&=\mathrm{i}\frac{1}{2}\big(\varphi^*(\partial_0+\mathrm{i}A_0)\varphi-\varphi(\partial_0-\mathrm{i}A_0)\varphi^*\big)-\frac{1}{2m}\big|(\partial_i+\mathrm{i}A_i)\varphi\big|^2\\
&\quad+\mu|\varphi|^2-\frac{V_0}{2}|\varphi|^4+\frac{1}{8\pi e^2}\left(\frac{1}{c}\pmb{E}^2-c\pmb{B}^2\right) \tag{3.7.5}
\end{align*}
其中 $c$ 是光速，而
\begin{equation*}
E_i=\partial_0A_i-\partial_iA_0=F_{0i},\qquad B_i=\epsilon_{ijk}\partial_jA_k=\frac{1}{2}\epsilon_{ijk}F_{jk}
\end{equation*}
分别是 $U(1)$ 规范理论的电场和磁场。

我们知道，具有整体 $U(1)$ 对称性的模型有一个守恒荷。借助规范场可以容易地证明这一点。规范化后的作用量是规范不变的：
\begin{equation*}
S(\varphi,A_\mu)=S(\mathrm{e}^{\mathrm{i}f(\pmb{x},t)}\varphi,A_\mu-\partial_\mu f)
\end{equation*}
设 $\varphi_c(\pmb{x},t)$ 是经典运动方程的一个解，那么
\begin{equation*}
S(\mathrm{e}^{\mathrm{i}f(\pmb{x},t)}\varphi_c,A_\mu)=S(\varphi_c,A_\mu)+O(f^2)
\end{equation*}
因此
\begin{align*}
S(\varphi_c,A_\mu)&=S(\varphi_c,A_\mu-\partial_\mu f)+O(f^2)\\
&=S(\varphi_c,A_\mu)+\int\mathrm{d}^d\pmb{x}\,\mathrm{d}t\,\partial_\mu f\,J^\mu(\varphi_c,A_\mu)+O(f^2)
\end{align*}
其中 $J^\mu$ 是流
\begin{equation*}
J^\mu(\varphi,A_\mu)\equiv-\partial_{A_\mu}L(\varphi,A_\mu)
\end{equation*}
我们看到，如果 $\varphi(\pmb{x},t)$ 满足经典运动方程，则对任意 $f$ 都有 $\int\mathrm{d}^d\pmb{x}\,\mathrm{d}t\,\partial_\mu f\,J^\mu(\varphi_c,A_\mu)=0$，于是流 $J^\mu(\varphi_c,A_\mu)$ 守恒：
\begin{equation*}
\partial_\mu J^\mu(\varphi_c,A_\mu)=\partial_t\rho+\partial_iJ^i=0
\end{equation*}
其中 $\rho=J^0$ 是密度，$J^i$ 是流。当然，上述结果对 $A_\mu$ 场为零的情形同样成立，我们有 $\partial_\mu J^\mu(\varphi_c)=0$，这正是中性系统的流守恒。

对我们的玻色子系统，求得守恒流为
\begin{gather*}
J^0=\rho=\varphi^*\varphi\\
J^i=-\frac{\mathrm{i}}{2m}\big[\varphi^*(\partial_i\varphi)-(\partial_i\varphi^*)\varphi\big]+A_i|\varphi|^2
\end{gather*}
有意思的是，我们玻色子系统中的流依赖于规范势。

\subsection*{问题 3.7.1}

考虑一个与规范场耦合的格点玻色子系统：
\begin{equation*}
L=\mathrm{i}\frac{1}{2}\sum_i\big(\varphi_i^*(\partial_0+\mathrm{i}A_0(i))\varphi_i-\varphi_i(\partial_0-\mathrm{i}A_0(i))\varphi_i^*\big)+\sum_{\langle ij\rangle}\big(t_{ij}\varphi_j^*\varphi_i\,\mathrm{e}^{-\mathrm{i}a_{ij}}+\mathrm{h.c.}\big)
\end{equation*}
其中求和遍及所有对 $\langle ij\rangle$，$a_{ij}$ 是规范场 $a_{ij}=\int_i^j\mathrm{d}\pmb{x}\cdot\pmb{A}$。
\begin{enumerate}
\item 证明 $L$ 在格点规范变换（lattice gauge transformation）
\begin{equation*}
\varphi_i\rightarrow\mathrm{e}^{\mathrm{i}f_i(t)}\varphi_i,\qquad A_0(i)\rightarrow A_0(i)-\partial_0f_i,\qquad a_{ij}\rightarrow a_{ij}-f_j+f_i
\end{equation*}
下不变。
\item 求 $\varphi_i(t)$ 的运动方程。
\item 计算密度的时间导数，即 $\partial_0\varphi_i^*\varphi_i$。证明可以引入一个定义在链上的流 $J_{ij}$，并恢复格点流守恒关系
\begin{equation*}
\partial_0\varphi_i^*\varphi_i+\sum_jJ_{ij}=0
\end{equation*}
\item 证明 $J_{ij}$ 可以表示为 $-\partial L/\partial a_{ij}$。
\end{enumerate}

\subsection{流关联函数与电磁响应}

\begin{keypoints}
\item 流守恒对流关联施加了约束。许多重要的物理量，如压缩率和电导率（conductivity），都由流关联决定。
\item 要得到正确的响应，重要的是按正确的次序取 $\pmb{k}\to0$ 与 $\omega\to0$ 这两个极限。
\end{keypoints}

现在让我们计算系统对外加规范势的响应。我们想看看规范势 $A_\mu$ 能产生多大的流 $J^\mu$。按下式引入 $j^\mu$：
\begin{equation*}
j^0=\rho,\qquad j^i=-\frac{\mathrm{i}}{2m}\big[\varphi^*(\partial_i\varphi)-(\partial_i\varphi^*)\varphi\big]
\end{equation*}
我们看到拉格朗日量具有形式
\begin{equation*}
L(\varphi,A_\mu)=L(\varphi)-A_0j^0-A_ij^i-\frac{1}{2m}(A^i)^2\rho
\end{equation*}
利用线性响应理论计算 $\langle j^\mu(\pmb{x},t)\rangle$ 到 $A_\mu$ 的领头阶，我们得到流 $J^\mu\equiv-\partial_{A_\mu}L$：
\begin{equation*}
\langle J^\mu(\pmb{x},t)\rangle=\langle j^\mu(\pmb{x},t)\rangle+(1-\delta^{\mu0})A^\mu\rho=\int\mathrm{d}^d\pmb{x}\,\mathrm{d}t\,\Pi^{\mu\nu}(\pmb{x},t;\pmb{x}',t')A_\nu(\pmb{x}',t')
\end{equation*}
其中响应函数为
\begin{align*}
\Pi^{00}(\pmb{x},t;\pmb{x}',t')&=-\mathrm{i}\Theta(t-t')\,\big\langle[\rho(\pmb{x},t),\rho(\pmb{x}',t')]\big\rangle\\
\Pi^{0i}(\pmb{x},t;\pmb{x}',t')&=-\mathrm{i}\Theta(t-t')\,\big\langle[\rho(\pmb{x},t),j^i(\pmb{x}',t')]\big\rangle\\
\Pi^{i0}(\pmb{x},t;\pmb{x}',t')&=-\mathrm{i}\Theta(t-t')\,\big\langle[j^i(\pmb{x},t),\rho(\pmb{x}',t')]\big\rangle\\
\Pi^{ij}(\pmb{x},t;\pmb{x}',t')&=-\mathrm{i}\Theta(t-t')\,\big\langle[j^i(\pmb{x},t),j^j(\pmb{x}',t')]\big\rangle+\delta^{ij}\delta(\pmb{x}-\pmb{x}')\delta(t-t')\frac{\langle\rho\rangle}{m}
\end{align*}
由于流依赖于 $A_i$，我们多出一个接触项（contact term）$\delta^{ij}\delta(\pmb{x}-\pmb{x}')\delta(t-t')\frac{\langle\rho\rangle}{m}$。如果我们引入关联函数
\begin{equation*}
\pi^{\mu\nu}(\pmb{x},t;\pmb{x}',t')=-\mathrm{i}\Theta(t-t')\,\big\langle[j^\mu(\pmb{x},t),j^\nu(\pmb{x}',t')]\big\rangle
\end{equation*}
则
\begin{equation*}
\Pi^{\mu\nu}=\pi^{\mu\nu}+\delta^{\mu\nu}(1-\delta^{0\mu})(1-\delta^{0\nu})\delta(\pmb{x}-\pmb{x}')\delta(t-t')\langle\rho\rangle \tag{3.7.6}
\end{equation*}
上述结果对零温和有限温都成立。我们注意到
\begin{equation*}
\big(\Pi^{\mu\nu}(\pmb{x},t;\pmb{x}',t')\big)^*=\Pi^{\nu\mu}(\pmb{x}',t';\pmb{x},t)
\end{equation*}
或者，在 $\omega$--$\pmb{k}$ 空间中，
\begin{equation*}
\big(\Pi^{\mu\nu}_{(k_\lambda)}\big)^*=\Pi^{\nu\mu}_{(-k_\lambda)}
\end{equation*}
其中 $k_0=\omega$。

由于流守恒，$\Pi^{\mu\nu}$ 的各分量并非完全独立。在 $\omega$--$\pmb{k}$ 空间中，我们有
\begin{equation*}
k_\mu\Pi^{\mu\nu}_{k_\lambda}=0
\end{equation*}
于是，完整的 $\Pi^{\mu\nu}$ 可以由 $\Pi^{ij}$ 按如下方式确定：
\begin{align*}
\Pi^{0i}_{(k_\lambda)}&=\big(\Pi^{i0}_{(-k_\lambda)}\big)^\dagger=-\frac{k_j}{\omega}\Pi^{ji}_{(k_\lambda)}\\
\Pi^{00}_{(k_\lambda)}&=-\frac{k_j}{\omega}\Pi^{0j}_{(k_\lambda)}=\frac{k_ik_j}{\omega^2}\Pi^{ij}_{(k_\lambda)} \tag{3.7.7}
\end{align*}
对于转动不变的系统，我们还可以把 $\Pi^{ij}$ 和 $\pi^{ij}$ 进一步分解成一个纵向分量（longitudinal component）$\Pi^{\parallel}_{(k_\lambda)}$ 和一个横向分量（transverse component）$\Pi^{\perp}_{(k_\lambda)}$：
\begin{align*}
\Pi^{ij}_{(k_\lambda)}&=\frac{k_ik_j}{\pmb{k}^2}\Pi^{\parallel}_{(k_\lambda)}+\Big(\delta_{ij}-\frac{k_ik_j}{\pmb{k}^2}\Big)\Pi^{\perp}_{(k_\lambda)}\\
\pi^{ij}_{(k_\lambda)}&=\frac{k_ik_j}{\pmb{k}^2}\pi^{\parallel}_{(k_\lambda)}+\Big(\delta_{ij}-\frac{k_ik_j}{\pmb{k}^2}\Big)\pi^{\perp}_{(k_\lambda)} \tag{3.7.8}
\end{align*}
式 (3.7.7) 现在可以写成
\begin{equation*}
\Pi^{0i}_{(k_\lambda)}=-k_i\Pi^{\parallel}_{(k_\lambda)},\qquad \Pi^{00}_{(k_\lambda)}=\frac{\pmb{k}^2}{\omega^2}\Pi^{\parallel}_{(k_\lambda)}
\end{equation*}
而式 (3.7.6) 可以写成
\begin{equation*}
\Pi^{\parallel}_{(k_\lambda)}=\pi^{\parallel}_{(k_\lambda)}+\frac{\langle\rho\rangle}{m},\qquad \Pi^{\perp}_{(k_\lambda)}=\pi^{\perp}_{(k_\lambda)}+\frac{\langle\rho\rangle}{m}, \tag{3.7.9}
\end{equation*}
响应函数 $\Pi^{\mu\nu}$ 与许多重要的物理量相关。让我们以金属为例。在 $\omega\to0$ 极限下，我们有
\begin{equation*}
\delta\rho(\pmb{k})=\Pi^{00}_{(0,\pmb{k})}A_0(\pmb{k})
\end{equation*}
注意 $A_0(\pmb{x})$ 正是外加势，而 $-\Pi^{00}_{(0,\pmb{k})}$ 就是（波矢 $\pmb{k}$ 处的）压缩率：
\begin{equation*}
\chi(\pmb{k})=-\Pi^{00}_{(0,\pmb{k})}
\end{equation*}
通常，我们预期 $\Pi^{00}_{(0,\pmb{k})}$ 对所有 $\pmb{k}$ 都有限。于是，在小 $\omega$ 以及 $\omega\ll\pmb{k}$ 的极限下（注意这里我们先取 $\omega\to0$），我们有
\begin{equation*}
\Pi^{\parallel}_{(k_\lambda)}=-\chi(\pmb{k})\frac{\omega^2}{\pmb{k}^2}
\end{equation*}
由式 (3.7.9) 可见，在 $\omega\to0$ 极限下，$\pi^{\parallel}_{(k_\lambda)}$ 必须恰好抵消 $\frac{\langle\rho\rangle}{m}$，才能使 $\Pi^{\parallel}_{(k_\lambda)}$ 像 $\omega^2$ 那样趋于零。

现在让 $\pmb{k}$ 先趋于零。在 $|\pmb{k}|\ll\omega$ 极限下，$-\mathrm{i}\omega A_{i,(\omega,\pmb{k})}$ 是一个近乎均匀的电场，预期它会产生一个方向由 $A_i$ 给出、波矢由 $\pmb{k}$ 给出的流：
\begin{equation*}
J^i(\omega)=\lim_{\pmb{k}\to0}\frac{\Pi^{ij}}{-\mathrm{i}\omega}(-\mathrm{i}\omega)A_{j,(\omega,\pmb{k})}
\end{equation*}
由式 (3.7.8) 我们看到，只有当 $\omega\gg|\pmb{k}|$ 极限下 $\Pi^{\parallel}_{(k_\lambda)}=\Pi^{\perp}_{(k_\lambda)}$ 时，$\pmb{k}\to0$ 的极限才存在。假设情形正是如此，那么我们得到电导率
\begin{equation*}
\sigma(\omega)=\frac{\Pi^{\parallel}_{(\omega,0)}}{-\mathrm{i}\omega}=\frac{\Pi^{\perp}_{(\omega,0)}}{-\mathrm{i}\omega}
\end{equation*}
电导率的实部
\begin{equation*}
\mathrm{Re}\,\sigma(\omega)=-\mathrm{Im}\frac{\Pi^{\parallel}_{(\omega,0)}}{\omega}=-\mathrm{Im}\frac{\Pi^{\perp}_{(\omega,0)}}{\omega}
\end{equation*}
对应于耗散。电导率的虚部给出介电常数
\begin{equation*}
\epsilon(\omega)=-\frac{\mathrm{Im}\,\sigma(\omega)}{\omega}=\mathrm{Re}\frac{\Pi^{\parallel}_{(\omega,0)}}{\omega^2}=\mathrm{Re}\frac{\Pi^{\perp}_{(\omega,0)}}{\omega^2}
\end{equation*}
于是，在 $\omega\gg|\pmb{k}|$ 极限下，我们可以写成
\begin{equation*}
\Pi^{\parallel}_{(\omega,0)}=\Pi^{\perp}_{(\omega,0)}=\mathrm{i}\omega\,\mathrm{Re}\,\sigma(\omega)+\omega^2\epsilon(\omega)
\end{equation*}
我们注意到，极化矢量 $\pmb{P}$ 满足 $\partial_{\pmb{x}}\cdot\pmb{P}=-\delta\rho$，或者（假定 $A_i=0$）
\begin{equation*}
\mathrm{i}k_iP^i=-\delta\rho=-\Pi^{00}A_0=-\frac{\pmb{k}^2}{\omega^2}\Pi^{\parallel}_{(k_\lambda)}A_0=\mathrm{i}\frac{\Pi^{\parallel}_{(k_\lambda)}}{\omega^2}k_iE_i
\end{equation*}
因此
\begin{equation*}
P^i=\frac{\Pi^{\parallel}_{(k_\lambda)}}{\omega^2}E_i
\end{equation*}
我们又看到 $\Pi^{\parallel}_{(k_\lambda)}/\omega^2$ 是介电常数。

我们已经考虑了 $\Pi^{\parallel}$ 和 $\Pi^{\perp}$ 的 $\omega\gg|\pmb{k}|$ 极限，以及 $\Pi^{\parallel}$ 的 $\omega\ll|\pmb{k}|$ 极限。那么 $\Pi^{\perp}$ 的 $\omega\ll|\pmb{k}|$ 极限是什么呢？一个自然的猜测是它对应于磁化率。磁矩密度 $\pmb{M}$ 满足 $\partial_{\pmb{x}}\times\pmb{M}=-\pmb{j}$。于是（假定 $A_0=0$），我们有
\begin{align*}
\mathrm{i}\epsilon^{ijk}k_jM_k&=-j^i=-\Pi^{ij}A_j=-(\delta^{ij}\pmb{k}^2-k^ik^j)A_j\frac{\Pi^{\perp}_{(k_\lambda)}}{\pmb{k}^2}\\
&=+\epsilon^{ij'k'}k_{j'}\epsilon^{k'i'j}k_{i'}A_j\frac{\Pi^{\perp}_{(k_\lambda)}}{\pmb{k}^2}=-\mathrm{i}\epsilon^{ijk}k_jB_k\frac{\Pi^{\perp}_{(k_\lambda)}}{\pmb{k}^2}
\end{align*}
因此
\begin{equation*}
M_i=-\frac{\Pi^{\perp}_{(k_\lambda)}}{k^2}B_i
\end{equation*}
而 $-\Pi^{\perp}_{(k_\lambda)}/k^2$ 就是磁化率。

现在让我们计算玻色子超流体的 $\Pi^{\mu\nu}$。我们从带规范场 (3.7.3) 的玻色子拉格朗日量出发，积掉振幅涨落，得到一个带规范场的 XY 模型。保留到 $(\partial_\mu\theta,A_\mu)$ 的二次阶，我们有
\begin{equation*}
L=\frac{\chi}{2}\big((\partial_0\theta+A_0)^2-v^2(\partial_i\theta+A_i)^2\big) \tag{3.7.10}
\end{equation*}
由此我们看到
\begin{equation*}
j^0=-\chi\partial_0\theta,\qquad j^i=\chi v^2\partial_i\theta
\end{equation*}
因此
\begin{align*}
\pi^{00}&=\chi^2(-\mathrm{i}\omega)(\mathrm{i}\omega)\big\langle\theta_{(\omega,\pmb{k})}\theta_{(-\omega,-\pmb{k})}\big\rangle=\chi\frac{\omega^2}{\omega^2-v^2\pmb{k}^2+\mathrm{i}0^+\mathrm{sgn}(\omega)}\\
\pi^{0i}&=\pi^{i0}=\chi^2(-\mathrm{i}\omega)(-\mathrm{i}k_i)\big\langle\theta_{(\omega,\pmb{k})}\theta_{(-\omega,-\pmb{k})}\big\rangle=-\chi v^2\frac{\omega k_i}{\omega^2-v^2\pmb{k}^2+\mathrm{i}0^+\mathrm{sgn}(\omega)}\\
\pi^{ij}&=\chi^2(\mathrm{i}k_i)(-\mathrm{i}k_j)\big\langle\theta_{(\omega,\pmb{k})}\theta_{(-\omega,-\pmb{k})}\big\rangle=\chi v^4\frac{k_ik_j}{\omega^2-v^2\pmb{k}^2+\mathrm{i}0^+\mathrm{sgn}(\omega)}
\end{align*}
注意，$\mathrm{i}0^+\mathrm{sgn}(\omega)$ 的选择给出的正是响应函数。总的 $\Pi^{\mu\nu}$ 为
\begin{align*}
\Pi^{00}&=\chi\left(\frac{\omega^2}{\omega^2-v^2\pmb{k}^2+\mathrm{i}0^+\mathrm{sgn}(\omega)}-1\right)\\
\Pi^{0i}&=\Pi^{i0}=-\chi v^2\frac{\omega k_i}{\omega^2-v^2\pmb{k}^2+\mathrm{i}0^+\mathrm{sgn}(\omega)}\\
\Pi^{ij}&=\chi v^2\left(\delta_{ij}+v^2\frac{k_ik_j}{\omega^2-v^2\pmb{k}^2+\mathrm{i}0^+\mathrm{sgn}(\omega)}\right)
\end{align*}
而
\begin{gather*}
\Pi^{\parallel}=\chi v^2\left(\frac{v^2k^2}{\omega^2-v^2\pmb{k}^2+\mathrm{i}0^+\mathrm{sgn}(\omega)}+1\right)\\
\Pi^{\perp}=\chi v^2
\end{gather*}
压缩率 $-\Pi^{00}_{(0,\pmb{k})}$ 有限且等于 $\chi$。磁化率 $-\Pi^{\perp}/c\pmb{k}^2$ 在 $\pmb{k}\to0$ 时发散。电导率的实部
\begin{equation*}
\mathrm{Re}\,\sigma(\omega)=-\mathrm{Im}\frac{\Pi^{\parallel}_{(\omega,0)}}{\omega}=\mathrm{Im}\frac{\chi v^2}{\omega+\mathrm{i}0^+}=\pi\frac{\rho}{m}\delta(\omega)
\end{equation*}
在有限频率处为零。

如果我们选取库仑规范（Coulomb gauge）$\partial_{\pmb{x}}\cdot\pmb{A}=0$，那么由 $J^i=\Pi^{ij}A_j$，我们发现流与\emph{规范势}之间的一个简单关系：
\begin{equation*}
\pmb{J}=\Pi^{\perp}\pmb{A}=\frac{\rho}{m}\pmb{A} \tag{3.7.11}
\end{equation*}
这就是著名的伦敦方程（London equation）。超导体的许多新奇性质——如持续电流（persistent current）、迈斯纳效应等——都由它负责。

\subsection{超流性与有限温度效应}

\begin{keypoints}
\item 激发会削弱超流，但无法消灭它。这导致了超流性。
\item 超流只能靠涡旋的隧穿来消灭。
\item 超流的临界速度。
\end{keypoints}

最后，我们准备讨论相互作用玻色子对称性破缺相的超流性质（Landau, 1941）。首先考虑自由空间中的玻色子系统。该玻色子系统在伽利略变换（Galileo transformation）下不变。让我们假定对称性破缺基态之上的激发具有谱 $\epsilon(\pmb{k})$，并忽略激发之间的相互作用。后一假定在低温、激发稀薄时成立。

考察（静止的）基态之上的一个单激发 $(\epsilon,\pmb{k})$。系统的总能量与总动量分别为 $E_{\mathrm{ground}}+\epsilon$ 与 $\pmb{P}=\pmb{k}$。如果我们以速度 $\pmb{v}$ 推动（boost）系统，则系统的总能量与总动量将分别变为 $E=E_{\mathrm{ground}}+\epsilon+\frac{1}{2}Nm\pmb{v}^2+\pmb{v}\cdot\pmb{k}$ 与 $\pmb{P}=\pmb{k}+Nm\pmb{v}$。与被推动基态的能量和动量（即 $E=E_{\mathrm{ground}}+\frac{1}{2}Nm\pmb{v}^2$ 与 $\pmb{P}=Nm\pmb{v}$）相比，我们看到被推动基态之上的激发具有新谱
\begin{equation*}
\epsilon_{\pmb{v}}(\pmb{k})=\epsilon(\pmb{k})+\pmb{v}\cdot\pmb{k}
\end{equation*}
在被推动的超流体中，动量为 $\pmb{k}$ 的激发的占据数由 $n_B(\epsilon(\pmb{k}))$ 给出，其中 $n_B(\epsilon)=1/(\mathrm{e}^{\epsilon/T}-1)$。被推动超流体的能量与动量分别可写为
\begin{gather*}
\tilde{E}=E_{\mathrm{ground}}+\sum_{\pmb{k}}\epsilon_{\pmb{v}}(\pmb{k})n_B(\epsilon(\pmb{k}))+\frac{1}{2}Nm\pmb{v}^2\\
\tilde{\pmb{P}}=\sum_{\pmb{k}}\pmb{k}\,n_B(\epsilon(\pmb{k}))+Nm\pmb{v}=Nm\pmb{v}
\end{gather*}
在平衡态中，激发的占据数应为 $n_B(\epsilon_{\pmb{v}}(\pmb{k}))$。我们看到，被推动的超流体并不处于平衡态。如果我们让系统通过重新分配激发而弛豫到实验室系中的平衡态，那么能量与动量将分别变为
\begin{gather*}
E_{\pmb{v}}=E_{\mathrm{ground}}+\sum_{\pmb{k}}\epsilon_{\pmb{v}}(\pmb{k})n_B(\epsilon_{\pmb{v}}(\pmb{k}))+\frac{1}{2}Nm\pmb{v}^2\\
\pmb{P}_{\pmb{v}}=\sum_{\pmb{k}}\pmb{k}\,n_B(\epsilon_{\pmb{v}}(\pmb{k}))+Nm\pmb{v}
\end{gather*}
对于小的 $\pmb{v}$，上式可重写为\footnote{$\pmb{P}_{\pmb{v}}$ 的结果很容易看出。为得到 $E_{\pmb{v}}$，我们注意到
\begin{align*}
E_{\pmb{v}}&=E_{\mathrm{ground}}+\sum_{\pmb{k}}\epsilon(\pmb{k})n_B(\epsilon(\pmb{k}))+\sum_{\pmb{k}}(\pmb{v}\cdot\pmb{k})^2\big[n_B'(\epsilon(\pmb{k}))+\tfrac{1}{2}\epsilon(\pmb{k})n_B''(\epsilon(\pmb{k}))\big]+\frac{1}{2}Nm\pmb{v}^2\\
&=E_0-\frac{1}{2}\mathcal{V}m\rho_n\pmb{v}^2+\frac{1}{2d}\pmb{v}^2\sum_{\pmb{k}}\pmb{k}^2\frac{\mathrm{d}}{\mathrm{d}\epsilon(\pmb{k})}\big[\epsilon(\pmb{k})n_B'(\epsilon(\pmb{k}))\big]+\frac{1}{2}Nm\pmb{v}^2
\end{align*}}
\begin{align*}
E_{\pmb{v}}&=E_0+\frac{1}{2}\big[Nm-\mathcal{V}(\rho_nm+\delta\rho_nm)\big]\pmb{v}^2\\
\pmb{P}_{\pmb{v}}&=(Nm-\mathcal{V}\rho_nm)\pmb{v} \tag{3.7.12}
\end{align*}
其中
\begin{equation*}
\rho_n=-\frac{1}{md}\int\frac{\mathrm{d}^d\pmb{k}}{(2\pi)^d}\,\pmb{k}^2n_B'(\epsilon(\pmb{k})) \tag{3.7.13}
\end{equation*}
\begin{equation*}
\delta\rho_n=-\frac{1}{md}\int\frac{\mathrm{d}^d\pmb{k}}{(2\pi)^d}\,\pmb{k}^2\frac{\mathrm{d}}{\mathrm{d}\epsilon(\pmb{k})}\big(\epsilon(\pmb{k})n_B'(\epsilon(\pmb{k}))\big) \tag{3.7.14}
\end{equation*}
如果对小 $\pmb{k}$ 有 $\epsilon(\pmb{k})\propto|\pmb{k}|$，则对小 $T$ 有 $\rho_n\propto T^{d+1}$。

这里我们看到对称性破缺相的一个惊人性质：在我们让激发弛豫到平衡态之后，系统的总动量 $\pmb{P}_{\pmb{v}}$ 并不弛豫到零。平衡态中玻色子永远运动下去。正是这一性质让对称性破缺相得到了“超流体”这个名字。为了更好地理解这一现象，我们回顾一下：在对称性破缺相中，玻色子场可以写成
\begin{equation*}
\varphi=\varphi_0+\delta\varphi
\end{equation*}
其中 $\varphi_0$ 描写凝聚体，$\delta\varphi$ 描写凝聚体之上的激发。设玻色子处于线度为 $L$、带周期边界条件的盒子中。当我们推动系统时，我们把玻色子边界条件从 $\varphi(L)=\varphi(0)$ 扭转到 $\varphi(L)=\mathrm{e}^{\mathrm{i}mvL}\varphi(0)$。于是，在这一推动下，$\varphi$ 变为 $\mathrm{e}^{\mathrm{i}mvx}\varphi$。在被推动的超流体中，凝聚体被扭转：$\varphi_0\to\mathrm{e}^{\mathrm{i}mvx}\varphi_0$（见图 3.15(a)）。由于周期边界条件，$mvL$ 按 $2\pi\times$ 整数量子化。现在很清楚：由于 $|\varphi_0|$ 固定，我们不可能在不撕裂凝聚体的情况下解除扭转。解除凝聚体的扭转需要把凝聚体 $\varphi_0$ 压到零，这要耗费巨大的能量。更确切地说，解除扭转需要产生一个涡旋，并让这个涡旋一路移动、横穿整个样品（见图 3.15(b)）。因此，解除凝聚体的扭转是一个具有有限能量势垒的隧穿过程。按照这幅图像，我们看到：如果忽略隧穿过程，运动的超流体尽管能量较高，也不能弛豫回基态。我们还学到，超流体并不能真正地永远流动。隧穿会使流动弛豫到零，不过这可能需要非常非常长的时间。超流动性的关键在于无穷大系统中涡旋的无穷大能量代价，而这一无穷大的能量代价来自式 (3.4.2) 中有限的相位劲度（phase rigidity）$\eta\neq0$。

然而，对激发来说，情况就大不相同了。这里 $\delta\varphi$ 在零附近涨落，可以扭转并轻易地改变自己的相位。因此，涨落（即激发）能够通过其与环境的相互作用（这种相互作用可以非常弱）轻易地弛豫到平衡态。

上述讨论暗示了一幅两流体（two-fluid）图像。玻色子超流体包含两个组分：与凝聚体相联系的超流组分，以及与激发相联系的正常流体组分。当玻色子超流体流过管道时，只有超流组分能够无摩擦地流过。摩擦太大时，正常流体组分无法流动。（在稳态下管道内没有压差，正常流体组分便完全无法流动。）当然，超流的速度不能太大。如果 $\pmb{v}$ 太大，实验室系中的激发能量 $\epsilon_{\pmb{v}}(\pmb{k})$ 可能变成负值，这意味着不稳定以及无摩擦流动的终结。因此，临界速度为（见图 3.1）
\begin{equation*}
v_c=\mathrm{Min}\big(\epsilon(\pmb{k})/|\pmb{k}|\big)
\end{equation*}

% ==================================================================
% 新增术语（ch3_f 新增，供术语表追加）：
%   minimal coupling 最小耦合；contact term 接触项；
%   conductivity 电导率；longitudinal/transverse component 纵向分量 / 横向分量；
%   Coulomb gauge 库仑规范；polarization vector 极化矢量；
%   magnetic moment density 磁矩密度；London equation 伦敦方程；
%   Meissner effect 迈斯纳效应；persistent current 持续电流；
%   lattice gauge transformation 格点规范变换；
%   Galileo transformation / Galileo invariance 伽利略变换 / 伽利略不变性；
%   boost 推动；two-fluid picture 两流体图像；
%   superfluid component / normal fluid component 超流组分 / 正常流体组分；
%   phase rigidity 相位劲度；energy barrier 能量势垒；
%   laboratory frame 实验室系；field strength 场强
% ==================================================================
